\chapter{어텐션 메커니즘}
\label{ch:attention}

이 장에서는 트랜스포머의 핵심인 \keyterm{어텐션(attention)} 메커니즘을 구현한다. Scaled Dot-Product Attention에서 시작해 Multi-Head Attention까지 단계적으로 살펴본다.

\section{직관: 어텐션이란 무엇인가}

문장 ``The cat sat on the mat because \uline{it} was tired''에서 ``it''이 무엇을 가리키는가? 사람은 쉽게 ``cat''임을 안다. 컴퓨터는 어떻게 알 수 있을까? 어텐션은 이 문제를 푸는 메커니즘이다. 각 단어가 다른 모든 단어에 ``얼마나 주의를 기울일지''를 학습하여, ``it''이 ``cat''에 높은 가중치를 주도록 만드는 것이다.

수식으로, 어텐션은 주어진 쿼리(query) $Q$에 대해, 모든 키(key) $K$와의 유사도를 계산하고, 그 유사도로 값(value) $V$를 가중 합산하는 연산이다.

\begin{formula}[어텐션의 일반적 정의]
\[
\text{Attention}(Q, K, V) = \text{softmax}\left(\frac{Q \cdot K^T}{\sqrt{d_k}}\right) \cdot V
\]
\begin{itemize}
  \item $Q$: 쿼리 (지금 보고 있는 토큰의 표현)
  \item $K$: 키 (다른 토큰들의 ``라벨'')
  \item $V$: 값 (다른 토큰들의 ``내용'')
  \item $d_k$: 키의 차원. $\sqrt{d_k}$로 나누는 이유는 dot product가 커져 softmax가 포화되는 것을 방지하기 위함.
\end{itemize}
\end{formula}

\subsection{직관적 비유: 도서관 검색}

도서관에서 책을 찾는 상황을 상상하자. 당신이 원하는 정보(쿼리 $Q$)를 가지고 도서관에 간다. 각 책에는 제목과 키워드(키 $K$)가 있고, 내용(값 $V$)이 있다. 도서관 시스템은 당신의 쿼리를 각 책의 키와 비교하여 유사도를 계산하고, 가장 유사한 책의 내용을 가중 합산하여 답을 준다.

어텐션은 이 과정을 미분 가능하게 만든 것이다. ``유사도''를 dot product로, ``가중 합산''을 softmax 가중치로 수행한다.

\section{Q, K, V는 어디서 오는가}

트랜스포머에서 $Q, K, V$는 모두 같은 입력 $x$에서 선형 변환으로 만들어진다.

\[
Q = x W_Q, \quad K = x W_K, \quad V = x W_V
\]

여기서 $W_Q, W_K, W_V$는 학습 가능한 행렬. 이렇게 하면 같은 입력에 대해 ``질문'', ``라벨'', ``내용'' 세 가지 관점을 학습할 수 있다.

\begin{infobox}[왜 세 개의 다른 행렬인가?]
같은 $x$를 세 번 사용하지 않고 $W_Q, W_K, W_V$로 변환하는 이유는 ``역할 분담'' 때문이다. 쿼리는 ``무엇을 찾을지'', 키는 ``무엇이 있는지'', 값은 ``실제 내용''을 표현해야 한다. 같은 벡터로 세 역할을 모두 하기 어렵다. 선형 변환을 통해 각 역할에 맞는 ``표상 공간''으로 투영하는 것이다.
\end{infobox}

\section{스케일링: 왜 $\sqrt{d_k}$로 나누는가?}

$Q \cdot K^T$는 $d_k$차원 벡터들의 내적이다. $d_k$가 크면 내적 값이 커진다. $Q, K$의 원소가 평균 0, 분산 1이라고 하면, 내적 $\sum_{i=1}^{d_k} Q_i K_i$의 분산은 $d_k$다.

내적 값이 크면 softmax 출력이 원-핫에 가까워진다 (한 값이 1, 나머지 0). 이렇게 되면 그래디언트가 거의 0이 되어 학습이 멈춘다. $\sqrt{d_k}$로 나누면 분산이 1이 되어 softmax가 부드럽게 분포한다.

\begin{formula}[스케일링의 효과]
\[
\text{softmax}\left(\frac{Q K^T}{\sqrt{d_k}}\right) \quad \text{vs} \quad \text{softmax}(Q K^T)
\]
스케일링 없으면 $d_k = 64$일 때 내적 값이 $\pm 64$까지 갈 수 있어 softmax가 포화. 스케일링하면 $\pm 8$ 범위로 줄어 부드러운 분포.
\end{formula}

\section{Scaled Dot-Product Attention 구현}

가장 기본 형태부터 시작하자. PyTorch 2.0의 \code{F.scaled\_dot\_product\_attention}(SDPA)을 사용하면 메모리 효율적이고 빠르다.

\begin{lstlisting}[language=Python]
import torch
import torch.nn.functional as F

class ScaledDotProductAttention(nn.Module):
    def __init__(self, dropout: float = 0.0):
        super().__init__()
        self.dropout_p = dropout

    def forward(self, q, k, v, mask=None):
        # q, k, v: [batch, n_heads, seq, d_head]
        # mask: [batch, 1, seq, seq] 또는 [seq, seq]
        out = F.scaled_dot_product_attention(
            q, k, v, attn_mask=mask,
            dropout_p=self.dropout_p if self.training else 0.0,
        )
        return out
\end{lstlisting}

SDPA는 C++ 레벨에서 최적화되어 있어, 파이썬 루프로 softmax를 직접 구현하는 것보다 수십 배 빠르다. 또한 Flash Attention 같은 최신 알고리즘을 자동으로 활용한다.

\begin{notebox}[Flash Attention]
Flash Attention (Dao et al. 2022)은 어텐션 계산을 GPU 메모리 계층 구조에 맞게 재설계하여, 메모리 접근을 줄이고 속도를 2$\sim$4배 향상시킨 알고리즘이다. PyTorch 2.0의 SDPA는 자동으로 Flash Attention을 활성화한다. 우리가 직접 구현할 필요 없이 \code{F.scaled\_dot\_product\_attention} 한 줄로 혜택을 받는다.
\end{notebox}

\section{인과 마스크 (Causal Mask)}

자기회귀 언어 모델에서는 ``미래''를 보면 안 된다. 토큰 $t$를 예측할 때 토큰 $t+1, t+2, \dots$를 보면 안 된다는 뜻이다. 이를 위해 어텐션 점수 행렬의 상삼각(주대각선 위)을 $-\infty$로 마스킹하여 softmax 후 0이 되도록 한다.

\begin{lstlisting}[language=Python]
def make_causal_mask(seq_len: int, device="cpu") -> torch.Tensor:
    """하삼각 마스크. 미래 토큰을 -inf로 마스킹."""
    mask = torch.full((seq_len, seq_len), float("-inf"), device=device)
    mask = torch.triu(mask, diagonal=1)  # 주대각선 위를 -inf로
    return mask
\end{lstlisting}

\begin{figure}[htbp]
\centering
\begin{tikzpicture}[
  scale=0.7,
  every node/.style={font=\small\sffamily}
]
  % 5x5 마스크 시각화
  \foreach \i in {0, 1, 2, 3, 4} {
    \foreach \j in {0, 1, 2, 3, 4} {
      \pgfmathsetmacro{\val}{ifthenelse(\j <= \i, "0", "-inf")}
      \pgfmathsetmacro{\col}{ifthenelse(\j <= \i, "inkblue!10", "inkred!15")}
      \node[draw=inkgray, fill=\col, minimum size=0.8cm] at (\j*0.8, -\i*0.8) {\val};
    }
  }
  % 라벨
  \node[inkgray, anchor=east] at (-0.5, 0) {token 0};
  \node[inkgray, anchor=east] at (-0.5, -3.2) {token 4};
  \node[inkgray, anchor=south] at (0, 0.5) {token 0};
  \node[inkgray, anchor=south] at (3.2, 0.5) {token 4};
\end{tikzpicture}
\caption{인과 마스크 (5$\times$5). 파란색 = 어텐션 허용(0), 빨간색 = 미래 토큰 차단($-\infty$). 행 $i$는 토큰 $i$가 열 $j$ 토큰에 얼마나 주의를 기울일지 나타냄.}
\label{fig:causal-mask}
\end{figure}

마스크를 어텐션 점수에 더하면, 마스킹된 위치의 점수는 $-\infty$가 되고, softmax 후에는 0이 된다. 즉 미래 토큰의 $V$가 출력에 기여하지 않게 된다.

\subsection{왜 $-\infty$인가?}

softmax의 수식 $\text{softmax}(x_i) = e^{x_i} / \sum_j e^{x_j}$에서 $x_i = -\infty$면 $e^{-\infty} = 0$이 된다. 따라서 마스킹된 위치의 가중치가 정확히 0이 된다. 단순히 0을 더하는 것과 다르다: 0을 더하면 원래 점수가 그대로 남아 softmax 후 여전히 양수가 된다.

\section{Multi-Head Attention}

단일 어텐션은 하나의 ``관점''만 학습할 수 있다. 하지만 언어에는 다양한 관계가 있다: 문법적 관계(주어-동사), 의미적 관계(동의어, 반의어), 담화 관계(지시어 참조 해결) 등. \keyterm{멀티헤드 어텐션(Multi-Head Attention)}은 이를 여러 헤드로 분할하여 병렬 처리한다.

\subsection{아이디어}

$d_{model}$ 차원을 $h$개의 헤드로 분할하여, 각 헤드가 $d_{head} = d_{model}/h$ 차원에서 독립적으로 어텐션을 수행한다. 예를 들어 $d_{model}=512, h=8$이면 $d_{head}=64$.

\begin{formula}[Multi-Head Attention]
\[
\text{MHA}(x) = \text{Concat}(\text{head}_1, \dots, \text{head}_h) W_O
\]
\[
\text{head}_i = \text{Attention}(Q W_Q^i, K W_K^i, V W_V^i)
\]
각 헤드의 출력을 이어붙인 후 $W_O$로 선형 변환.
\end{formula}

\subsection{왜 헤드를 나누는가?}

$d_{model} = 512$인 단일 어텐션과 $d_{head} = 64, h = 8$인 멀티헤드 어텐션은 총 파라미터 수가 같다. 하지만 멀티헤드는 각 헤드가 다른 패턴을 학습할 수 있다. 예를 들어:
\begin{itemize}
  \item 헤드 1: 주어-동사 관계
  \item 헤드 2: 형용사-명사 수식
  \item 헤드 3: 지시어 참조
  \item 헤드 4: 문장 경계
\end{itemize}

실제로 어텐션 헤드를 시각화해보면 (예: BertViz), 서로 다른 헤드가 확실히 다른 패턴에 주목하는 것을 볼 수 있다.

\begin{figure}[htbp]
\centering
\begin{tikzpicture}[
  every node/.style={font=\small\sffamily},
  box/.style={draw=inkblue, fill=inkblue!8, rounded corners=2pt, minimum width=2.2cm, minimum height=0.6cm},
  head/.style={draw=inkred, fill=inkred!10, rounded corners=2pt, minimum width=1.6cm, minimum height=0.5cm, font=\scriptsize\sffamily}
]
  \node[box] (in) at (0, 0) {input [batch, seq, d\_model]};
  \node[box, below=0.5cm of in] (qkv) {QKV proj $\to$ split};
  % 헤드들
  \node[head] (h1) at (-3, -2) {head 1};
  \node[head] (h2) at (-1, -2) {head 2};
  \node[inkgray, font=\scriptsize] at (1, -2) {$\cdots$};
  \node[head] (h3) at (3, -2) {head h};
  \node[box, below=0.6cm of $(h1)!0.5!(h3)$] (cat) {Concat $\to$ O proj};
  \node[box, below=0.4cm of cat] (out) {output [batch, seq, d\_model]};

  \draw[-{Stealth[length=4pt]}, inkblue] (in) -- (qkv);
  \draw[-{Stealth[length=4pt]}, inkblue] (qkv) -- (h1);
  \draw[-{Stealth[length=4pt]}, inkblue] (qkv) -- (h2);
  \draw[-{Stealth[length=4pt]}, inkblue] (qkv) -- (h3);
  \draw[-{Stealth[length=4pt]}, inkblue] (h1) -- (cat);
  \draw[-{Stealth[length=4pt]}, inkblue] (h2) -- (cat);
  \draw[-{Stealth[length=4pt]}, inkblue] (h3) -- (cat);
  \draw[-{Stealth[length=4pt]}, inkblue] (cat) -- (out);
\end{tikzpicture}
\caption{Multi-Head Attention 구조. 입력을 $h$개 헤드로 분할하여 병렬 어텐션 수행 후 합침.}
\label{fig:mha}
\end{figure}

\subsection{구현}

효율성을 위해 $Q, K, V$를 한 번에 계산한다. 입력 $x$에 대해 $3 d_{model}$ 차원으로 프로젝션한 뒤, 헤드별로 분할한다.

\begin{lstlisting}[language=Python]
class MultiHeadAttention(nn.Module):
    def __init__(self, config: ModelConfig):
        super().__init__()
        self.d_model = config.d_model
        self.n_heads = config.n_heads
        self.d_head = config.d_model // config.n_heads
        # Q, K, V를 한 번에 계산 (결합 프로젝션)
        self.qkv_proj = nn.Linear(config.d_model, 3 * config.d_model, bias=config.use_bias)
        self.o_proj = nn.Linear(config.d_model, config.d_model, bias=config.use_bias)
        self.attention = ScaledDotProductAttention(dropout=config.dropout)
        self.dropout = nn.Dropout(config.dropout)

    def forward(self, x, mask=None):
        batch, seq, _ = x.shape
        # [batch, seq, 3*d_model] -> [batch, seq, 3, n_heads, d_head]
        qkv = self.qkv_proj(x).view(batch, seq, 3, self.n_heads, self.d_head)
        # [3, batch, n_heads, seq, d_head]
        qkv = qkv.permute(2, 0, 3, 1, 4)
        q, k, v = qkv[0], qkv[1], qkv[2]

        if mask is not None and mask.dim() == 2:
            mask = mask.unsqueeze(0).unsqueeze(0)  # [1, 1, seq, seq]

        # 어텐션: [batch, n_heads, seq, d_head]
        attn_out = self.attention(q, k, v, mask=mask)
        # 헤드 합치기: [batch, seq, d_model]
        attn_out = attn_out.transpose(1, 2).contiguous().view(batch, seq, self.d_model)
        return self.dropout(self.o_proj(attn_out))
\end{lstlisting}

\begin{notebox}[왜 QKV를 한 번에 계산하는가?]
세 개의 \code{nn.Linear}를 따로 쓰는 대신 \code{3 * d\_model} 출력을 가진 하나의 \code{nn.Linear}를 쓰는 이유는 효율성이다. GPU는 큰 행렬 연산을 한 번에 수행하는 것에 최적화되어 있다. 세 번의 작은 연산보다 한 번의 큰 연산이 더 빠르다. 분할은 \code{view + permute}로 메모리 재구성 없이 수행된다.
\end{notebox}

\section{인과성 테스트}

어텐션 모듈의 가장 중요한 테스트는 ``인과 마스크가 실제로 미래 토큰을 차단하는가''다. 다음 테스트는 이를 검증한다.

\begin{lstlisting}[language=Python]
def test_attention_is_causal(small_config):
    """위치 i의 출력이 위치 i+1 이후에 영향을 받지 않아야 함."""
    attn = MultiHeadAttention(small_config)
    attn.eval()
    x = torch.randn(1, 6, small_config.d_model)
    mask = make_causal_mask(6)
    out1 = attn(x, mask=mask)
    # x의 마지막 위치를 바꿔도 앞 위치의 출력은 변하지 않아야 함
    x2 = x.clone()
    x2[:, -1] = torch.randn(1, small_config.d_model)
    out2 = attn(x2, mask=mask)
    # 앞 5개 위치는 동일해야 함
    assert torch.allclose(out1[:, :5], out2[:, :5], atol=1e-5)
\end{lstlisting}

이 테스트가 통과하면 인과 마스크가 올바르게 작동하고 있음을 확신할 수 있다. 만약 실패한다면 마스크가 뒤집혀 있거나 차원이 잘못된 것이다.

\section{Self-Attention vs Cross-Attention}

이 장에서 구현한 것은 \keyterm{자기 어텐션(self-attention)}으로, $Q, K, V$가 모두 같은 시퀀스에서 온다. 번역 모델의 디코더에서는 \keyterm{교차 어텐션(cross-attention)}도 사용하는데, 이는 $Q$는 디코더에서, $K, V$는 인코더에서 온다. GPT는 디코더 only 구조라 자기 어텐션만 사용한다.

\subsection{교차 어텐션은 어떻게 다른가?}

교차 어텐션에서 $Q = x_{dec} W_Q$, $K = x_{enc} W_K$, $V = x_{enc} W_V$다. 디코더의 각 위치가 인코더 출력의 모든 위치를 참조한다. 번역에서 디코더가 ``번역할 문장''을 생성할 때 원문의 해당 단어를 찾는 과정이다.

GPT는 이 교차 어텐션이 없다. 오직 자기 어텐션만으로 다음 토큰을 예측한다. 이것이 GPT를 ``디코터 only'' 모델이라 부르는 이유다.

\section{어텐션 계산량 분석}

어텐션의 계산량은 $O(L^2 \cdot d)$다. $L$은 시퀀스 길이. 시퀀스가 길어지면 제곱으로 증가한다.

\begin{center}
\begin{tabularx}{\textwidth}{lXX}
\toprule
\textbf{시퀀스 길이} & \textbf{어텐션 행렬 크기} & \textbf{FP16 메모리 (단일 헤드)} \\
\midrule
512 & $512 \times 512$ & 0.5 MB \\
2048 & $2048 \times 2048$ & 8 MB \\
8192 & $8192 \times 8192$ & 128 MB \\
32768 & $32768 \times 32768$ & 2 GB \\
\bottomrule
\end{tabularx}
\end{center}

이것이 긴 문맥 처리가 어려운 이유다. 2권 7장에서 KV cache와 PagedAttention으로 이 비용을 줄이는 방법을 다룬다.

\section{실습: 어텐션 가중치 시각화}

학습된 모델의 어텐션 가중치를 시각화하면 모델이 어떤 패턴에 주목하는지 볼 수 있다.

\begin{lstlisting}[language=Python]
import matplotlib.pyplot as plt
import math
from makellm.model.attention import make_causal_mask

# 어텐션 가중치 수동 계산 및 시각화 (SDPA는 가중치를 내부적으로만 관리하므로 디버그 시 직접 계산)
model.eval()
with torch.no_grad():
    x = torch.randn(1, 10, model.config.d_model) # [batch=1, seq=10, d_model]
    x_norm = model.blocks[-1].ln1(x)
    qkv = model.blocks[-1].attn.qkv_proj(x_norm)
    batch, seq, _ = x.shape
    qkv = qkv.view(batch, seq, 3, model.config.n_heads, model.config.d_head).permute(2, 0, 3, 1, 4)
    q, k = qkv[0], qkv[1]  # [1, n_heads, seq, d_head]
    
    # 0번째 헤드의 가중치 계산: (Q @ K^T) / sqrt(d_head)
    scores = (q[0, 0] @ k[0, 0].T) / math.sqrt(model.config.d_head)
    mask = make_causal_mask(seq)
    attn_weights = torch.softmax(scores + mask, dim=-1)  # [seq, seq]

plt.imshow(attn_weights.numpy(), cmap='viridis')
plt.colorbar()
plt.xlabel('key position')
plt.ylabel('query position')
plt.title('Attention weights (head 0, last layer)')
\end{lstlisting}

\section{요약}

\begin{itemize}
  \item 어텐션은 쿼리-키 유사도로 값을 가중 합산하는 연산이다.
  \item $\sqrt{d_k}$로 스케일링하여 softmax 포화를 방지한다.
  \item 인과 마스크로 미래 토큰을 차단하여 자기회귀 모델이 된다.
  \item 멀티헤드는 $d_{model}$을 $h$개 헤드로 분할하여 다양한 관계를 학습.
  \item QKV를 한 번에 계산하여 GPU 효율을 극대화.
  \item PyTorch 2.0의 SDPA를 활용하면 최적화된 어텐션을 한 줄로 사용 가능.
  \item 어텐션 계산량은 $O(L^2)$로 긴 시퀀스에서 병목.
\end{itemize}

\section{연습 문제}

\begin{enumerate}
  \item $\sqrt{d_k}$ 스케일링을 빼면 학습이 어떻게 달라지는가? $d_k = 64$일 때 softmax 출력의 분포를 계산해보라.
  \item 인과 마스크를 씌우지 않으면 어떤 일이 일어나는가? 학습 시와 추론 시 각각 설명하라.
  \item 헤드 수 $h$를 1로 하면 (단일 헤드) 멀티헤드와 어떤 차이가 있는가?
  \item 헤드 수 $h$를 $d_{model}$과 같게 하면 (즉 $d_{head} = 1$) 어떤 일이 일어나는가?
  \item 어텐션 행렬 $L \times L$의 메모리가 $L = 100,000$일 때 얼마인지 계산하라 (FP16 기준).
\end{enumerate}

다음 장에서는 이 어텐션을 포함한 트랜스포머 블록을 조립한다. 어텐션 + FFN + 잔차 + LayerNorm의 조합이 트랜스포머를 만든다.
